% ===================================================================
%  HOTO Na 13 — Otal. --- Gidan cin abinci. --- Gidan kofi.
%
%  Traduction, faite en 2026, en haoussa d'aujourd'hui, sur l'ido de
%  texto/io. Regles de la colonne : voir 10-hoto-01.tex. Une seule
%  numerotation.
%
%  NEUVIEME REFUS DE LA COLONNE, ET C'EST LE MEME QUE LE JAVANAIS A
%  FAIT SUR LA MEME PLANCHE. Le javanais avait « dhakon » et
%  « macanan » pour les jeux du cafe et ne les a pas ecrits ; le
%  haoussa a « dara », qui se joue dans des trous creuses en terre
%  avec des cailloux et qui est le jeu de toute la savane. La planche
%  grave un damier, un jeu de dominos, un trictrac et des cartes sur
%  des tables de marbre. On n'ecrit donc pas dara.
%
%  MAIS LES ECHECS (78) SONT ACCEPTES SANS UNE HESITATION, ET C'EST
%  LE CONTRAIRE EXACT. « sadaranji » vient de l'arabe شطرنج, qui vient
%  du persan chatrang, qui vient du sanskrit chaturanga — et le mot
%  francais « echecs » vient du meme endroit par la meme route. Ce
%  n'est pas un emprunt qu'on ferait aujourd'hui : c'est le meme mot
%  que le francais, arrive ici par le Sahara au lieu de l'Espagne.
%  Deux langues qui se rencontrent sur un mot venu d'Inde.
%
%  ET LE TRICTRAC (81) EST « wasan tebur », LE JEU DE LA TABLE — ce
%  qui n'est pas une trouvaille mais une constatation. L'arabe dit
%  طاولة, la table ; l'italien dit tavole ; l'anglais disait tables
%  avant de dire backgammon ; le vieux francais disait le jeu des
%  tables, et trictrac ne nomme que le bruit. Le haoussa a exactement
%  le meme droit de l'appeler par son meuble, et il ne fait que
%  rejoindre quatre langues qui l'ont fait avant lui.
%
%  DEUX MOTS POUR DEUX AUBERGES, ET NEUF TABLEAUX ENTRE ELLES. Celle
%  du tableau 7 etait « masauki », le lieu ou l'on est loge, avec son
%  enseigne qui grincait au vent du nord et son aubergiste qui fumait
%  sur le seuil. Celle-ci est « otal » : six etages, un ascenseur, un
%  telephone, une caisse, un registre des voyageurs, un groom a
%  bicyclette. La colonne n'a pas hesite une seconde entre les deux
%  mots, et c'est le meilleur exemple qu'elle donne de ce que « on
%  modernise la langue, jamais les choses » veut dire dans les deux
%  sens : on ne modernise pas non plus une auberge de village en lui
%  donnant le mot d'un palace.
%
%  ET LES TRIPES (une des entrees du menu) SONT « tumbi », QUI ETAIT
%  L'ESTOMAC AU TABLEAU 2. Huitieme paire d'homonymes de la colonne,
%  et la plus courte : entre l'organe et le plat, il n'y a que la
%  cuisson. Le haoussa ne voit pas la difference et il a raison.
%
%  ENFIN LE MENU EST FRANCAIS D'UN BOUT A L'AUTRE et la colonne le
%  garde tel quel — crevettes, tripes a la mode de Caen, gigot,
%  epinards a la creme, Saint-Emilion. Les epinards (alayyaho) sont
%  le seul plat de la liste qui pousse ici, et c'est aussi le seul
%  mot haoussa de tout l'alinea.
% ===================================================================

%%K t13-noto-1 noto
\VUnotes{143.2mm}{%
\textit{\VUgras{(*) Game da cikakken bayanin ɗaki, duba hoto na 6.}}%
}

%%K t13-apar-1 apar
\VUsaut{3.0mm}
\VUtitre{15.0pt}{60}{HOTO Na 13}
\VUsaut{1.6mm}
\VUpk{10.2pt}{40}{Otal. --- Gidan cin abinci.}
\VUsaut{2.0mm}
\VUpk{10.2pt}{40}{Gidan kofi.}
\VUsaut{3.4mm}

%%K t13-01-1 p
1. --- Da na iso Royan jiya da yamma, na sauka a « \textit{Otal na
Babban Teku} » wanda wani aboki ya ba ni shawarar zuwa.

%%K t13-01-2 p
An ba ni kyakkyawan \VUgras{ɗaki} \textsuperscript{(2)}, a kan \VUgras{bene}
\textsuperscript{(1)} na biyu inda na yi barci sosai (*).

%%K t13-02-1 p
2. --- Yau da safe, ina fitowa daga ɗakina, inda \VUgras{mai aikin
gida} \textsuperscript{(3)} ta kawo ƙaramin abincin safe, sai na shiga
\VUgras{ɗakin shan taba} \textsuperscript{(5)} wanda ake amfani da shi a
matsayin \VUgras{ɗakin karatu} \textsuperscript{(4)} domin karanta
\VUgras{jaridu} \textsuperscript{(6)}, ana shan \VUgras{taba}
\textsuperscript{(8)}. Bayan na karanta, na sauka da \VUgras{lif}
\textsuperscript{(9)} na kuma je yawo cikin garin.

%%K t13-03-1 p
3. --- Tun da na manta in tambayi \VUgras{ma’aikacin ofis}
\textsuperscript{(12)} na otal ɗin ko a wane lokaci ake ba da abinci a kan
\VUgras{teburin gamayya} \textsuperscript{(11)}, sai na komo da wuri na kuma
yi amfani da hakan in je in yi wanka a \VUgras{ɗakin wanka}
\textsuperscript{(10)} na otal ɗin. Sa’an nan na sake zuwa \VUgras{ofishin
kuɗi} \textsuperscript{(15)} domin in yi wa ɗaya daga cikin abokaina waya.
Amma ana amfani da \VUgras{tarho} \textsuperscript{(13)} ; da ta ga na
kunce, \VUgras{ma’aikaciyar kuɗi} \textsuperscript{(14)}, wadda take rubuta
\VUgras{lissafi} \textsuperscript{(17)} a kan \VUgras{littafin matafiya}
\textsuperscript{(16)}, sai ta roƙi \VUgras{mai gadin ƙofa}
\textsuperscript{(18)} ya aiki \VUgras{ɗan aike} \textsuperscript{(19)} ya yi
mini saƙo \VUgras{a kan keke} \textsuperscript{(20)}.

%%K t13-04-1 p
4. --- Na gode mata na kuma fita domin in ba da ɗan kyauta ga
\VUgras{ɗan ɗako} \textsuperscript{(24)} wanda ya kawo daga tashar a kan
\VUgras{amalankensa} \textsuperscript{(25)} \VUgras{akwatin hulata}
\textsuperscript{(26)}, da \VUgras{jakar kayana} \textsuperscript{(27)} da
\VUgras{jakar tafiyata} \textsuperscript{(28)}. \VUgras{Mai kai wasiƙa}
\textsuperscript{(21)}, wanda ya iso yanzu, ya ba ni \VUgras{wasiƙa}
\textsuperscript{(22)}.

%%K t13-05-1 p
5. --- Na ɗan daɗe a bakin ƙofa, kusa da \VUgras{akwatin wasiƙa}
\textsuperscript{(23)}, ina kallon zirga-zirga a kan hanya.
\VUgras{Matuƙin} \textsuperscript{(30)} \VUgras{babbar motar haya}
\textsuperscript{(29)} yana ɗora a kan abin hawansa \VUgras{akwatin}
\textsuperscript{(31)} wani \VUgras{matafiyi} \textsuperscript{(32)} wanda
\VUgras{mai otal} \textsuperscript{(33)} yake raka. Wani matashin \VUgras{mai
waya} \textsuperscript{(35)} ya kawo \VUgras{waya} \textsuperscript{(34)} ba tare
da gaggawa ba domin wani abokin cinikin otal ; wani \VUgras{matafiyi}
\textsuperscript{(36)} da \VUgras{kyamararsa} \textsuperscript{(38)} a kafaɗa,
yana duba \VUgras{littafin jagorarsa} \textsuperscript{(37)}.

%%K t13-tit-1 sub
\VUsaut{4.0mm}
\VUpk{10.2pt}{40}{Lokacin abincin rana.}
\VUsaut{3.0mm}

%%K t13-06-1 p
6. --- A gaban shingen, wani \VUgras{ɗan aike} \textsuperscript{(39)} marar
damuwa yana gyangyaɗi tsakanin \VUgras{ƙugiyar ɗaukarsa}
\textsuperscript{(40)} da \VUgras{akwatin gogen takalminsa}
\textsuperscript{(41)}, alhali wata matashiyar \VUgras{mai wanki}
\textsuperscript{(42)}, wadda take ɗauke da \VUgras{kwando}
\textsuperscript{(43)} na tufafi, take dariya saboda fuskar tsanani ta
\VUgras{ma’aikacin tebur} \textsuperscript{(44)} wanda yake tsaye kusa da
ƙofar.

%%K t13-07-1 p
7. --- Ina ganin \VUgras{mai girki} \textsuperscript{(45)}, wanda ya fito
daga \VUgras{ɗakin girkinsa} \textsuperscript{(47)} da yake a \VUgras{benen
ƙarƙashin ƙasa} \textsuperscript{(48)} domin ya aiki \VUgras{mai wanke
kwanuka} \textsuperscript{(46)} ya yi saƙo, sai na tuna cewa lokacin abincin
rana ya kusa, sai na shiga gidan cin abinci, ina ƙetare farfajiya inda
\VUgras{matuƙin mota} \textsuperscript{(50)} ya ajiye \VUgras{motarsa}
\textsuperscript{(49)}, alhali \VUgras{makanike} \textsuperscript{(52)} yake
gyara, bisa ga umarnin \VUgras{mai keke} \textsuperscript{(53)},
\VUgras{keke mai ƙafa uku} \textsuperscript{(51)} wanda ya sami hatsari
yanzu.

%%K t13-08-1 p
8. --- Na tambayi matashin \VUgras{ɗan aike} \textsuperscript{(54)} wanda
yake ɗauke da \VUgras{gilashin giya} \textsuperscript{(55)}, ko teburin
gamayya yana ƙarƙashin baranda, kuma saboda amsarsa ta tabbatarwa, sai
na sami wuri a kan ɗaya daga cikin teburoci na kuma sa aka ba ni abinci
mai kyau ƙwarai.

%%K t13-noto-2 noto
\VUnotes{143.2mm}{%
\textit{\VUgras{(*) Da taimakon jerin kwanuka da aka shigar a ƙamus,
malami zai iya sauya jerin abincin da yake ƙasa cikin sauƙi.}}%
}

%%K t13-tit-2 sub
\VUsaut{4.0mm}
\VUpk{10.2pt}{40}{Abincin rana mai kyau ƙwarai.}
\VUsaut{3.0mm}

%%K t13-09-1 p
9. --- Ma’aikacin ya kawo mini jerin abinci (*) na rana da jerin giya.
Na zaɓa na kuma yi oda na \VUgras{jatan lande} a matsayin
\VUgras{abin buɗe ci} tare da \VUgras{man shanu,} kuma, a matsayin
abincin farko, \VUgras{tumbi bisa ga salon Caen.} Ban ci \VUgras{kifi}
ba, sai dai na nemi rabo na \VUgras{cinyar} ɗan tunkiya, kuma, a
matsayin \VUgras{kwanon ganyayyaki, alayyaho} da kirim. Sa’an nan, tun
da babu ɗaya daga cikin \VUgras{abinci na tsakiya} da ya ba ni sha’awa,
sai na sa aka kawo mini kayan zaki iri iri, \VUgras{cuku,} da
\VUgras{ƴaƴan itace} da \VUgras{biskit.} Na ƙara kuma \VUgras{giyar}
Saint-Emilion mai kyau ƙwarai, kuma na gamsu ƙwarai da abincina sai na
bar teburin bayan na ba \VUgras{ma’aikaci} \textsuperscript{(57)}
kyakkyawan \VUgras{kyauta} \textsuperscript{(59)}.

%%K t13-10-1 p
10. --- Ina ganin ni a shirye domin tafiya, sai \VUgras{manaja}
\textsuperscript{(60)} ya ba ni shawara in je baranda domin in yaba wa
kallon \VUgras{Babban Teku} \textsuperscript{(62)} mai kyau ƙwarai. Can nesa
sai muka fara ganin \VUgras{ƙananan jiragen} kamun kifi
\textsuperscript{(63)}, da \VUgras{jiragen ruwa masu zanen iska}
\textsuperscript{(64)} da babban \VUgras{jirgin fasinja}
\textsuperscript{(65)} wanda baƙin siffarsa yake bayyana a kan farin
\VUgras{gajimare} \textsuperscript{(67)} na \VUgras{iyakar ido}
\textsuperscript{(66)}. Iska mai laushi ta sa \VUgras{tuta}
\textsuperscript{(70)} da aka kafa a saman \VUgras{alama}
\textsuperscript{(69)} ta \VUgras{zauren} \textsuperscript{(68)} ta yi motsi.

%%K t13-tit-3 sub
\VUsaut{3.0mm}
\VUpk{10.2pt}{40}{Nishaɗi.}
\VUsaut{3.0mm}

%%K t13-11-1 p
11. --- Kallon ya ba da sha’awa ƙwarai, har na yanke shawara in hau
kan dandalin gidan kofi gobe ina ɗauke da \VUgras{ƙaramin madubin gani}
\textsuperscript{(71)} ko \VUgras{madubin gani mai nisa}
\textsuperscript{(72)}.

%%K t13-12-1 p
12. --- Ina barin baranda, na nufi gidan kofi na otal ɗin domin in sha
\VUgras{abin sha} \textsuperscript{(73)}, ina wasan \VUgras{biliyad}
\textsuperscript{(75)}. Ba tare da tsayawa ba na ƙetare zauren kofi inda
masu sha da yawa suke wasan \VUgras{sadaranji} \textsuperscript{(78)}, da
\VUgras{wasan faifai} \textsuperscript{(79)}, da \VUgras{domino}
\textsuperscript{(80)}, da \VUgras{wasan tebur} \textsuperscript{(81)} ko da
\VUgras{kati} \textsuperscript{(82)}, na kuma hau \VUgras{ɗakin biliyad}
\textsuperscript{(74)} kai tsaye.

%%K t13-13-1 p
13. --- Lokacin da wasan ya ƙare, na sauka zuwa benen ƙasa na kuma
nemi wurin ma’aikaci \VUgras{ambulan} \textsuperscript{(84)}, da
\VUgras{takarda} \textsuperscript{(83)}, da \VUgras{kan sarki}
\textsuperscript{(85)}, da \VUgras{alƙalami} \textsuperscript{(87)} da
\VUgras{tawada} \textsuperscript{(86)} domin in rubuta wasiƙa.

%%K t13-13-2 p
Sa’an nan na kira \VUgras{matuƙi} \textsuperscript{(92)} wanda
\VUgras{abin hawansa} \textsuperscript{(88)} ya tsaya gaban gidan kofi. Na
tambaye shi farashi bisa ga sa’o’i ko na tafiya ɗaya, kuma, lokacin da
muka yarda a kan farashin, sai ya buɗe mini \VUgras{ƙofar abin hawa}
\textsuperscript{(89)} ya zaunar da ni. Ya sake hawa \VUgras{mazauninsa}
\textsuperscript{(91)}, ya tada dawakin da sauri da \VUgras{ƙaramin bulala}
\textsuperscript{(93)} sai muka tafi yawo mai kwarjini.

\VUsaut{6.0mm}
\VUfilet{22mm}
